% ═══════════════════════════════════════════════════════════════
% ARBEITSBLATT: El tiempo y los lugares
% Spanisch Klasse 8 - Unidad 6, DS 1
% ═══════════════════════════════════════════════════════════════

\documentclass[10pt, a4paper]{article}

\usepackage[utf8]{inputenc}
\usepackage[T1]{fontenc}
\usepackage[ngerman]{babel}

% --- SCHRIFTEN ---
\usepackage{palatino}
\usepackage{helvet}
\newcommand{\sanstext}[1]{{\fontfamily{phv}\selectfont #1}}

\usepackage{microtype}
\usepackage[margin=1.4cm, top=1.6cm, bottom=1.3cm]{geometry}
\usepackage{enumitem}
\usepackage{tabularx}
\usepackage{booktabs}
\usepackage{array}
\usepackage{multicol}
\usepackage{tikz}
\usetikzlibrary{calc}

\usepackage{fancyhdr}
\usepackage{lastpage}
\usepackage{needspace}

\usepackage{xcolor}
\usepackage{tcolorbox}
\tcbuselibrary{skins}

% ═══════════════════════════════════════════════════════════════
% FARBEN
% ═══════════════════════════════════════════════════════════════

\definecolor{kopfzeile}{HTML}{1A1A1A}
\definecolor{akzent}{HTML}{C41E3A}
\definecolor{hellgrau}{HTML}{F2F2F2}
\definecolor{mittelgrau}{HTML}{777777}
\definecolor{dunkelgrau}{HTML}{444444}

% ═══════════════════════════════════════════════════════════════
% VARIABLEN
% ═══════════════════════════════════════════════════════════════

\newcommand{\thema}{El tiempo y los lugares}
\newcommand{\klasse}{8}
\newcommand{\maxpunkte}{28}

% ═══════════════════════════════════════════════════════════════
% KOPF- UND FUSSZEILE
% ═══════════════════════════════════════════════════════════════

\pagestyle{fancy}
\fancyhf{}
\renewcommand{\headrulewidth}{0pt}
\renewcommand{\footrulewidth}{0pt}
\cfoot{\sanstext{\footnotesize\textcolor{mittelgrau}{\thepage\,/\,\pageref{LastPage}}}}

% ═══════════════════════════════════════════════════════════════
% BOXEN
% ═══════════════════════════════════════════════════════════════

\newtcolorbox{grammatikbox}{
    colback=hellgrau,
    colframe=hellgrau,
    boxrule=0pt,
    arc=0pt,
    left=8pt, right=8pt, top=8pt, bottom=6pt,
    borderline north={1.5pt}{0pt}{dunkelgrau}
}

\newtcolorbox{vokabelkasten}{
    colback=white,
    colframe=mittelgrau,
    boxrule=0.5pt,
    arc=0pt,
    left=8pt, right=8pt, top=6pt, bottom=6pt
}

% ═══════════════════════════════════════════════════════════════
% BEFEHLE
% ═══════════════════════════════════════════════════════════════

\newcommand{\es}[1]{\textbf{#1}}
\newcommand{\de}[1]{{\small\textit{\textcolor{mittelgrau}{#1}}}}
\newcommand{\luecke}[1]{\rule{#1}{0.4pt}}
\newcommand{\punktebox}[1]{\hfill\sanstext{\small[\_\_/#1]}}

\newcommand{\aufgabe}[2]{%
    \needspace{4\baselineskip}%
    \vspace{12pt}%
    \noindent\sanstext{\textbf{#1}\quad#2}%
    \vspace{6pt}%
    \nopagebreak[4]%
}

\newlist{teil}{enumerate}{1}
\setlist[teil]{label=\alph*), leftmargin=1.8em, itemsep=4pt, parsep=0pt, topsep=3pt}

\setlength{\parindent}{0pt}
\setlength{\parskip}{4pt}
\setlength{\columnsep}{24pt}

% ═══════════════════════════════════════════════════════════════
% DOKUMENT
% ═══════════════════════════════════════════════════════════════

\begin{document}

% ═══════════════════════════════════════════════════════════════
% HEADER
% ═══════════════════════════════════════════════════════════════

\noindent
\begin{tikzpicture}[remember picture, overlay]
    \fill[kopfzeile] (current page.north west) rectangle +(\paperwidth, -1cm);
    \node[anchor=west, text=white, font=\sanstext\large\bfseries]
        at ([xshift=1.4cm, yshift=-0.5cm]current page.north west) {\thema};
    \node[anchor=east, text=white, font=\sanstext\small]
        at ([xshift=-1.4cm, yshift=-0.5cm]current page.north east)
        {Spanisch\,·\,Klasse \klasse\,·\,Arbeitsblatt};
\end{tikzpicture}

\vspace{0.5cm}

% --- Name/Datum ---
\noindent
\sanstext{\small\textbf{Name:}} \luecke{5cm} \hfill
\sanstext{\small\textbf{Datum:}} \luecke{2.5cm}

\vspace{8pt}

% ═══════════════════════════════════════════════════════════════
% VOKABELKASTEN (zum Ausfüllen)
% ═══════════════════════════════════════════════════════════════

\begin{vokabelkasten}
\textbf{Kasten 1: El tiempo} \de{-- Das Wetter} \hfill \textit{(Ergänze aus dem Lernblatt)}

\vspace{6pt}

\begin{multicols}{2}
\begin{tabular}{@{}p{4cm}p{4cm}@{}}
¿Qué tiempo hace? & \luecke{3.5cm} \\[0.6em]
Hace sol. & \luecke{3.5cm} \\[0.6em]
Hace calor. & \luecke{3.5cm} \\[0.6em]
Hace frío. & \luecke{3.5cm} \\[0.6em]
Hace viento. & \luecke{3.5cm} \\
\end{tabular}

\columnbreak

\begin{tabular}{@{}p{4cm}p{4cm}@{}}
Hace buen tiempo. & \luecke{3.5cm} \\[0.6em]
Hace mal tiempo. & \luecke{3.5cm} \\[0.6em]
Llueve. & \luecke{3.5cm} \\[0.6em]
Nieva. & \luecke{3.5cm} \\[0.6em]
Hay nubes. & \luecke{3.5cm} \\
\end{tabular}
\end{multicols}
\end{vokabelkasten}

\vspace{4pt}

\begin{vokabelkasten}
\textbf{Kasten 2: Los lugares} \de{-- Die Orte} \hfill \textit{(Ergänze aus dem Lernblatt)}

\vspace{6pt}

\begin{multicols}{3}
\begin{tabular}{@{}ll@{}}
el parque & \luecke{2cm} \\[0.5em]
la plaza & \luecke{2cm} \\[0.5em]
el cine & \luecke{2cm} \\[0.5em]
el teatro & \luecke{2cm} \\
\end{tabular}

\columnbreak

\begin{tabular}{@{}ll@{}}
la biblioteca & \luecke{2cm} \\[0.5em]
la iglesia & \luecke{2cm} \\[0.5em]
el supermercado & \luecke{2cm} \\[0.5em]
la panadería & \luecke{2cm} \\
\end{tabular}

\columnbreak

\begin{tabular}{@{}ll@{}}
la farmacia & \luecke{2cm} \\[0.5em]
el restaurante & \luecke{2cm} \\[0.5em]
la cafetería & \luecke{2cm} \\[0.5em]
el centro comercial & \luecke{2cm} \\
\end{tabular}
\end{multicols}
\end{vokabelkasten}

\vspace{6pt}

% ═══════════════════════════════════════════════════════════════
% AUFGABE 1: Partner-Interview
% ═══════════════════════════════════════════════════════════════

\aufgabe{1}{Partner-Interview: ¿Qué haces cuando...?} \punktebox{12}

Befrage deinen Partner / deine Partnerin und notiere die Antworten auf Spanisch.

\vspace{4pt}

\begin{tabularx}{\textwidth}{|p{4cm}|X|}
\hline
\textbf{Frage} & \textbf{Antwort deines Partners / deiner Partnerin} \\
\hline
\es{¿Qué haces cuando hace sol?} & \\[1.2em]
\de{Was machst du, wenn es sonnig ist?} & \\[0.8em]
\hline
\es{¿Qué haces cuando llueve?} & \\[1.2em]
\de{Was machst du, wenn es regnet?} & \\[0.8em]
\hline
\es{¿Qué haces cuando hace frío?} & \\[1.2em]
\de{Was machst du, wenn es kalt ist?} & \\[0.8em]
\hline
\es{¿Qué haces cuando hace calor?} & \\[1.2em]
\de{Was machst du, wenn es heiß ist?} & \\[0.8em]
\hline
\end{tabularx}

\vspace{4pt}
{\footnotesize\textit{Redemittel:} Cuando hace sol/frío/calor, voy al/a la... · Me quedo en casa. · Juego al...}

% ═══════════════════════════════════════════════════════════════
% AUFGABE 2: Wetter + Kleidung
% ═══════════════════════════════════════════════════════════════

\aufgabe{2}{¿Qué ropa llevas? Schreibe vollständige Sätze.} \punktebox{12}

\vspace{4pt}

\begin{teil}
    \item \es{¿Qué ropa llevas cuando hace mucho frío?}\\[0.3em]
    \luecke{14cm}

    \item \es{¿Qué ropa llevas cuando hace mucho calor?}\\[0.3em]
    \luecke{14cm}

    \item \es{¿Qué ropa llevas cuando llueve?}\\[0.3em]
    \luecke{14cm}

    \item \es{¿Qué ropa llevas cuando hace buen tiempo y vas al parque?}\\[0.3em]
    \luecke{14cm}
\end{teil}

\vspace{4pt}
{\footnotesize\textit{Redemittel:} Cuando hace..., llevo... · un abrigo, una bufanda, los guantes, el gorro, las botas, el paraguas, las gafas de sol, pantalones cortos, una camiseta...}

% ═══════════════════════════════════════════════════════════════
% AUFGABE 3: Beschreibe die Fotos
% ═══════════════════════════════════════════════════════════════

\aufgabe{3}{Beschreibe das Wetter auf den Fotos (Buch S. 102-103).} \punktebox{4}

\vspace{4pt}

\begin{multicols}{2}
\begin{tabular}{@{}ll@{}}
Foto 1: & \luecke{5cm} \\[0.8em]
Foto 2: & \luecke{5cm} \\[0.8em]
Foto 3: & \luecke{5cm} \\
\end{tabular}

\columnbreak

\begin{tabular}{@{}ll@{}}
Foto 4: & \luecke{5cm} \\[0.8em]
Foto 5: & \luecke{5cm} \\[0.8em]
Foto 6: & \luecke{5cm} \\
\end{tabular}
\end{multicols}

% ═══════════════════════════════════════════════════════════════
% FOOTER
% ═══════════════════════════════════════════════════════════════

\vfill

\noindent\rule{\textwidth}{0.5pt}

\vspace{4pt}

\hfill\sanstext{\textbf{Punkte:}} \luecke{1.2cm} / \maxpunkte

\end{document}
